\section{Introduction}\label{rev:sec:introduction}

Ramanujan's series for $1/\pi$ connect hypergeometric functions with
the arithmetic of elliptic curves. Their characteristic feature is a
coefficient linear in the summation index: a suitable combination of
a hypergeometric value and its first derivative becomes an algebraic
multiple of $1/\pi$. Modular parametrization explains many such
identities by evaluating at elliptic curves with complex multiplication
(CM). The converse asks whether the special value itself detects CM.
This distinction is central to a classification: constructing every
identity arising from CM does not exclude an isolated identity on a
non-CM fiber.

We study this converse for the four standard hypergeometric branches
of levels $1$--$4$, with primitive integer coefficients and an integer
denominator. There are three questions: whether the coefficient slope
is unique at a fixed parameter, whether every identity comes from CM,
and whether the CM parameters admit a finite arithmetic enumeration.
We prove uniqueness and carry out the CM classification unconditionally.
The implication from a single complex identity to CM remains
conditional on an explicit quadratic-period independence hypothesis.
The paper develops the period geometry of this implication and proves
unconditional arithmetic criteria and scarcity results around it.

\subsection{Prior work and the converse problem}
Ramanujan's 1914 paper \cite{Ramanujan1914} established the connection
between modular equations and rapidly convergent approximations to
$\pi$. Subsequent modular constructions organize these evaluations
through singular moduli and transformations of hypergeometric
functions. Chan--Chan--Liu \cite{CCL} give a general modular summation
identity and applications to Domb numbers; Chan--Cooper
\cite{ChanCooper} assemble rational analogues at levels $1$--$4$.
Chen--Glebov--Goenka \cite{CGG} derive formulas for noncompact arithmetic
triangle groups and identify their $36$ Clausen cases with the
level-$1$--$4$ list of Chan--Cooper. Their enumeration uses rational
uniformizer values at imaginary quadratic arguments. These works
establish the evaluations that appear in our standard tables and
provide the modular construction against which a converse can be tested.

Related constructions extend the range of coefficients, parameters,
and geometric realizations. Rogers \cite{RogersDomb} develops
hypergeometric transformations leading to further $1/\pi$ formulas,
and Almkvist--Guillera \cite{AG} study Ramanujan--Sato-like series.
Campbell--Cooper--Ye \cite{CCY} give quadratic irrational analogues and
explicit modular values. For compact Shimura curves, Yang \cite{Yang}
constructs Ramanujan-type identities whose normalization can involve
additional CM periods. These developments motivate the extensions in
our appendices. They also require care with the chosen branch, the
period target, and the convergence domain: these data belong to the
identity being classified.

The transcendence theory addresses a different part of the problem.
Chudnovsky's algebraic-independence results, in the precise forms
recorded by Waldschmidt \cite{Waldschmidt,WaldschmidtSurvey}, constrain
an elliptic period and its quasi-period strongly enough to prove
uniqueness of the coefficient slope. Nesterenko's theorem and the
period expressions for modular forms, presented in \cite{Fonseca},
give further independence statements. The quadratic-period independence
sufficient for the CM converse is stronger than the particular consequences of
these theorems used here. We isolate that relation and show how the
elliptic period conjecture and numerical modular-value conjectures
\cite{Fonseca,EterovicMSCD} would exclude it on non-CM fibers.
The results of Kreutz--Shen--Vial \cite{KSV} describe complete
de Rham--Betti comparison tensors; applying them requires the complete
comparison data, which one scalar special value does not supply.

Thus the $36$ standard evaluations are established identities. Our
purpose is to organize their arithmetic exhaustion within the CM
locus, identify the exact additional assertion needed for unrestricted
exhaustion, and relate that assertion to intrinsic period lines,
comparison tensors, and reduction at primes. The distinction between
these conclusions is retained throughout.

\subsection{Problem and principal results}

For $s\in\{1/6,1/4,1/3,1/2\}$ let
\[
 F_s(z)={}_3F_2\!\left(\begin{matrix}1/2,s,1-s\\1,1\end{matrix};z\right)
       =\sum_{n\ge0}c_s(n)z^n,
 \qquad c_s(n)=\frac{(1/2)_n(s)_n(1-s)_n}{(n!)^3},
 \qquad\thetaop=z\frac d{dz}.
\]
The pairs $(s_\ell,R_\ell)$ attached to levels $\ell=1,2,3,4$ are
\[
 (1/6,1728),\qquad(1/4,256),\qquad(1/3,108),\qquad(1/2,64).
\]
We use the power-series branch at zero and define
\begin{equation}\label{eq:series}
 \mathfrak S_\ell(A,B,C)
  =\sum_{n\ge0}R_\ell^nc_{s_\ell}(n)(A+Bn)C^{-n}
  =(A+B\thetaop)F_{s_\ell}(R_\ell/C),
\end{equation}
where the last expression always means the ordinarily convergent sum.

\begin{definition}
A primitive standard identity is an ordinarily convergent equality
\[
 \mathfrak S_\ell(A,B,C)=\frac\alpha\pi,
 \qquad A,B,C\in\Z,\quad C\ne0,\quad\gcd(A,B)=1,
 \quad\alpha\in\Qbar.
\]
Classes identify $(A,B,\alpha)$ with $(-A,-B,-\alpha)$, at the same
level and denominator. The multiplier is initially allowed to be zero.
\end{definition}

The primitive condition concerns the coefficient pair. Requiring only
$\gcd(A,B,C)=1$ would permit infinitely many multiples of the identity
$(A,B,C)=(1,6,256)$ at level four. Analytic continuation is likewise
excluded from the definition: at a convergence boundary, the value must
be the sum of the original series.

\begin{hypothesis}[Quadratic-period independence]\label{hyp:QPI}
Let $z=R_\ell/C\in[-1,1)\setminus\{0\}$, with $C\in\Z$, and put
$x=(1-\sqrt{1-z})/2$. For the algebraic elliptic curve $E_{s_\ell,x}$
of Section~\ref{sec:families}, use the analytic cycle normalized by
\[
 \omega=\frac{2\pi}{\sqrt6}
 {}_2F_1(s_\ell,1-s_\ell;1;x),
 \qquad h=\int_\gamma X\,\frac{dX}{y}.
\]
If this curve is non-CM, then $\pi,\omega^2,\omega h$ are linearly
independent over $\Qbar$.
\end{hypothesis}

\begin{theorem}[Conditional classification]\label{thm:main-conditional}
Assume Hypothesis~\ref{hyp:QPI}. Every primitive standard identity has
a CM elliptic fiber. There are exactly $36$ classes, comprising
$11,12,9,4$ classes at levels $1,2,3,4$, respectively, as listed in
Tables~\ref{tab:standard1}--\ref{tab:standard4}.
For each fixed level and denominator, the primitive coefficient pair
is unique up to simultaneous sign, and its multiplier is nonzero.

For the linear Euler companion $F_s(z)/(1-z)$, the complete list has
$35$ classes. It is obtained by the coefficient transport and ordinary
convergence criterion of Theorem~\ref{thm:euler}.
\end{theorem}
\begin{proof}
The elliptic realization and its first derivative identify the
coefficient space with de Rham cohomology. In the normalized basis,
\begin{equation}\label{rev:eq:leading-reduction}
 \pi\mathfrak S_\ell(A,B,C)
  =\frac{3}{2\pi}\bigl((A+Bu_s)\omega^2+Bv_s\omega h\bigr),
 \qquad u_s,v_s\in\Qbar,\quad v_s\ne0.
\end{equation}
Thus a nonzero coefficient pair gives a nontrivial quadratic-period
relation. Hypothesis~\ref{hyp:QPI} forces CM; the positive endpoint is
excluded unconditionally in Section~\ref{sec:endpoint}. CM reciprocity
and the finite modular correspondence then bound the class number and
give the tabulated parameters. The complementary CM eigenline supplies
the identity, and its equivariance supplies a rational coefficient
slope. These statements, with the exact finite enumeration, are proved
in Section~\ref{sec:enumeration}. Chudnovsky's theorem gives uniqueness
and nonvanishing in Section~\ref{sec:unique}. Finally, the invertible
Euler transport gives the companion classification.
\end{proof}

\begin{theorem}[Unconditional uniqueness]\label{thm:unconditional}
At every nonzero algebraic real $z\in[-1,1)$, there is at most one
$[A:B]\in\mathbf P^1(\Qbar)$ such that
$(A+B\thetaop)F_s(z)\in\Qbar/\pi$. A nonzero coefficient pair cannot
give zero. At $z=1$, no nonzero algebraic pair gives a convergent
identity in $\Qbar/\pi$.
\end{theorem}

The first two assertions are consequences of the identity-line theorem;
the endpoint uses CM period independence. In particular, the uniqueness
assertion in the leading theorem needs no conjecture.

\subsection{The period line and its arithmetic}
The organizing object is a line in the two-dimensional de Rham
cohomology of an elliptic curve. Clausen's identity realizes the
hypergeometric function as an elliptic period square. Its first
derivative is governed by the Gauss--Manin connection, which converts
the coefficient pair $(A,B)$ into a de Rham class. At an admissible
algebraic fiber this coefficient map is an isomorphism, and the
proposed identity takes the form
\[
 a\omega^2+b\omega h=\alpha\pi,
 \qquad a,b,\alpha\in\Qbar.
\]
Here $\omega$ is a holomorphic period and $h$ is the period of a
complementary algebraic class. Chudnovsky's theorem makes the resulting
coefficient line unique whenever it exists. At a CM fiber, the
endomorphism algebra supplies a canonical complement to the Hodge
line; the determinant comparison, or Legendre relation, gives the
required multiple of $1/\pi$. This construction explains both
uniqueness and CM existence in the same cohomological language.

For the converse, the natural identification
\[
 \operatorname{Sym}^2 H^1(E)\otimes\det H^1(E)^{-1}
       \simeq\operatorname{End}_0(H^1(E))
\]
places the quadratic relation in a comparison space of trace-zero
endomorphisms. A single scalar relation determines only part of that
comparison. The remaining coordinates, or suitable propagation to a
second cycle, would yield an endomorphism criterion for CM.
Quadratic-period independence gives an explicit numerical exclusion
sufficient for the original families. Once CM is established,
the modular correspondence bounds the degree of the elliptic
$j$-invariant by one or two. CM reciprocity \cite{CoxCM} and complete
class-number bounds \cite{Watkins} reduce the parameters to a finite
set. Descent of the complementary eigenline then gives the rational
slope before any coefficient identification is made. The linear Euler
companion is obtained by transporting this line; ordinary convergence
accounts for the change from $36$ standard classes to $35$ companions.

The same line also has an arithmetic realization. Finite
hypergeometric identities from \cite{BRSTT} relate weighted
truncations to the conjugate filtration in characteristic $p$.
Combined with the algebraization theorem of Tang \cite{Tang}, this
gives an unconditional characterization of the standard rows by
weighted vanishing at a set of primes of lower natural density
greater than $3/4$. At a single place of good ordinary reduction,
exact preservation of the Hodge line by the residue-field power of
crystalline Frobenius characterizes CM; the lifting step uses the
$p$-adic Lefschetz theorem of Maulik--Poonen \cite{MaulikPoonen}.
These criteria impose arithmetic information in addition to the
complex identity. They do not infer that information from its value.

Variation of the periods gives unconditional restrictions on possible
identities even without such arithmetic hypotheses. Product monodromy
and norm descent allow us to apply the Bombieri--Andr\'e specialization
theorem in the form stated in \cite{BaldiBinyaminiUrbanik2026}.
For bounded degree of the squared multiplier $\alpha^2$, the gaps
between distinct absolute denominators tend to infinity, jointly
across the four levels and both signs. On compact parameter intervals,
the counting theorem of Habegger--Pila \cite{HabeggerPilaCounting}
gives subpower bounds in complete arithmetic height windows. These
results restrict the distribution of possible identities; they leave
the existence of an isolated non-CM identity unresolved.

\subsection{Organization of the paper}
Section~\ref{sec:general-period-square} develops the intrinsic period
line, proves uniqueness, and treats ordinary convergence, including
the endpoints. Section~\ref{sec:hinge} formulates the CM implication
in terms of comparison tensors and modular derivatives, and proves
the implications from the stated numerical conjectures.
Section~\ref{sec:enumeration} gives CM descent, finite exhaustion,
and Euler transport. Section~\ref{sec:prime-tests} proves the
prime-density and one-prime criteria, and Section~\ref{sec:C-density}
establishes separation and counting.

The appendices supply the explicit elliptic models, the finite CM
certificate and all identity tables, and details of the integral
realizations. They also develop refinements of comparison and
specialization, finite transport, further elliptic branches, and
compact Shimura settings. The final appendices record the Pad\'e and
rational-certificate arguments, finite exclusions, and countertests.
Their role is to preserve the supporting arguments and delineate the
limits of stronger assertions without interrupting the main proofs.
